% TOP_PROBLEM: {"problemId": "problem.polynomial-time-solvability-of-condon-s-simple-stochastic-games", "canonicalTitle": "Polynomial-Time Solvability of Condon's Simple Stochastic Games", "releaseRank": 478, "primaryDomain": "theoretical_computer_science", "primaryDomainLabel": "Theoretical computer science", "displayStatus": "open", "releaseStatus": "open"}
% !TeX program = lualatex
% ATTEMPT_STATUS: unresolved
\documentclass[11pt]{article}
\usepackage[margin=25mm]{geometry}
\usepackage{fontspec}
\setmainfont{Latin Modern Roman}
\usepackage{amsmath,amssymb,mathtools}
\usepackage{unicode-math}
\setmathfont{Latin Modern Math}
\usepackage{xurl,hyperref}
\hypersetup{colorlinks=true,urlcolor=blue}
\setcounter{secnumdepth}{0}
\setlength{\parindent}{0pt}
\setlength{\parskip}{5pt}
\setlength{\emergencystretch}{3em}
\begin{document}
\section*{\texorpdfstring{Polynomial-Time Solvability of Condon's Simple Stochastic Games}{Polynomial-Time Solvability of Condon's Simple Stochastic Games}}\label{478}
\subsection{Definitions and mathematical statement}
A simple stochastic game is a finite directed graph whose nonterminal vertices belong to a maximizing player, a minimizing player, or chance. Players choose the outgoing edge at their own vertices; chance vertices use the given rational probabilities. Success means reaching a designated target, with nonreaching plays counted as failures. For initial vertex $s$, the value is
\[
 v(s)=\sup_\sigma\inf_\tau
 {\mathbb{P}}_s^{\sigma,\tau}(\text{eventually reach the target}).
\]
The question asks for a deterministic polynomial-time method to compute these values exactly, equivalently to decide whether $v(s)\ge q$ for a rational threshold $q$, in the standard finite rational-input convention. Binary decisions and transition encodings are part of the instance. An approximate answer with unspecified dependence on the precision or an algorithm exponential in input length does not settle the selected complexity question.
\par\smallskip\textit{Formulation and concept references:} \href{https://dblp.org/rec/journals/iandc/Condon92.html}{[S1]}.

\subsection{Short English statement}
Can the exact optimal winning probabilities of finite simple stochastic reachability games be computed in polynomial time?
\subsection{Research attempt}
\paragraph{Scope, conventions, and source check.}
This attempt by GPT-6 Astra Ultra was prepared on 27 September 2026.
The target includes games that can run forever, arbitrary explicitly
encoded rational chance probabilities, and the threshold relation
$v(s)\geq q$, including equality. Infinite nonreaching plays have payoff
zero. Make the target absorbing; make a dead end a failure self-loop.
Neither modification changes the event being measured. All nonterminal
vertices then have at least one outgoing transition.

Condon's original paper [E1] is the publication identified by the catalogue's
bibliographic link. Its abstract establishes the decision problem's
$\mathrm{NP}\cap\mathrm{coNP}$ membership, not a polynomial-time
algorithm. The primary algorithm paper [E2], Section 2, treats general
turn-based reachability games, states deterministic memoryless optimality,
and identifies the value as the least Bellman fixed point. It explicitly
discusses nonstopping components and the failure of arbitrary locally
optimal choices. The 2024 benchmark paper [E3] still introduces the
polynomial-time question as unresolved; its stopping-game experiments
are not a theorem for this full rational-input target. No general polynomial
algorithm was located in the consulted sources.

The attack first derives the least-fixed-point characterization with
nontermination handled explicitly, then solves the game exactly when both
stationary policies are fixed. A denominator bound shows that precision
and output size are not the central obstruction. Explicit slow examples
then defeat the attempt to obtain a polynomial algorithm by direct value
iteration. Finally a one-player linear program is examined to locate what
is lost when both strategic quantifiers are present.

\paragraph{Lemma 478.1: the least fixed point, without a stopping assumption.}
For a vector $x\in[0,1]^V$ with $x_t=1$ at the target, define
\[
 (Tx)_i=\begin{cases}
 1,&i=t,\\
 \max_{j\in\operatorname{Succ}(i)}x_j,&i\text{ belongs to Max},\\
 \min_{j\in\operatorname{Succ}(i)}x_j,&i\text{ belongs to Min},\\
 \sum_jp_{ij}x_j,&i\text{ is a chance vertex}.
 \end{cases}
\]
Failure terminals have value fixed at zero. Set $x^{(0)}_t=1$ and all
other initial coordinates to zero, and put $x^{(h+1)}=Tx^{(h)}$.
Then $x^{(h)}$ increases coordinatewise and
\begin{equation}\label{a478:least}
                         v=\lim_{h\to\infty}x^{(h)}
\end{equation}
is the least fixed point in this cube with the specified terminal values.

\emph{Proof.} Backward induction shows that $x^{(h)}_i$ is the value of
reaching the target within at most $h$ transitions: the last choice is a
maximum, minimum, or expectation according to the vertex type. Finite
action sets ensure each finite-horizon extremum is attained. The operator
is monotone, $x^{(0)}\leq Tx^{(0)}$, and all its iterates remain in
$[0,1]^V$. Thus they converge to a vector $z$. Finite maxima, minima,
and linear averages are continuous, giving $Tz=z$. Any other fixed point
$y$ in the cube dominates $x^{(0)}$ and therefore every iterate, so it
dominates $z$.

For each $h$, Max can use a finite-horizon optimal strategy for the first
$h$ moves and then continue arbitrarily. Its infinite-horizon success
probability against every Min strategy is at least $x^{(h)}_i$. Hence
$v(i)\geq z_i$. For the opposite inequality, choose at every Min vertex
a successor attaining the minimum in $Tz=z$, obtaining one stationary
Min strategy $\tau$. Under $\tau$ and any possibly randomized,
history-dependent Max strategy, $z$ of the current state is a bounded
supermartingale. At Max states every successor has $z_j\leq z_i$;
at the chosen Min transitions and at chance transitions its conditional
expectation equals $z_i$. After reaching the absorbing target it stays
equal to one. Thus the indicator of reaching the target by time $h$ is
at most $z$ of the state at time $h$, and its expectation is at most
$z_i$. Letting $h$ increase proves that the infinite-horizon reachability
probability is at most $z_i$. Therefore $v=z$. No interchange of an
infinite infimum with a probability limit is used.\hfill$\square$

The word ``least'' is indispensable. A single chance vertex with a
probability-one self-loop and no path to the target has equation $x=x$;
every $x\in[0,1]$ is a fixed point, but its reachability value is zero.
Similarly a Min vertex with actions ``stay here'' and ``go to target''
has equation $x=\min(x,1)$ and value zero. Solving the polynomial or
piecewise-linear equations without selecting the correct fixed point can
return a wrong value even in one state.

\paragraph{A polynomial graph computation of the zero-value region.}
Start with $W_0=\{t\}$ and repeatedly add a Max vertex with some
successor already present, a Min vertex with all successors present, or
a chance vertex with some positive-probability successor present.
Stop when no more vertices can be added, obtaining $W$. Then
\begin{equation}\label{a478:zero}
                          W=\{i:v(i)>0\}.
\end{equation}
To prove this, induction on the insertion round shows that every added
state has positive value for some finite horizon: at a Min vertex a
minimum of finitely many positive numbers remains positive, and at a
chance vertex one positive-probability positive term suffices. Conversely,
outside $W$, every Max successor stays outside, every positive chance
successor stays outside, and Min has at least one successor outside.
Choosing those Min successors keeps the play away from the target with
probability one. This proves \eqref{a478:zero}. Each round adds a vertex
or terminates, so the computation uses polynomial time in the explicit
graph size.

This qualitative computation is useful but does not find the positive
values. A region with values arbitrarily close to one and a region with
values exponentially small in the input length are both simply marked
``positive''. Replacing the quantitative question by this attractor test
would discard the rational threshold in the original target.

\paragraph{Lemma 478.2: exact evaluation after fixing both policies.}
Fix deterministic stationary policies $\sigma,\tau$, yielding a finite
Markov chain. Its target-hitting probabilities can be computed exactly in
polynomial bit complexity, including when the chain is not absorbing with
probability one.

\emph{Proof.} In the graph of positive-probability transitions, find the
states $R$ with a directed path to the target. Every state outside $R$
has hitting probability zero. Let $U=R\setminus\{t\}$, of size $r$,
and let $Q$ be the transition submatrix within $U$. Let $b_i$ be the
one-step probability of hitting $t$. The desired coordinates solve
\begin{equation}\label{a478:linear}
                             (I-Q)z=b.
\end{equation}
Transitions to states outside $R$ contribute zero rather than disappearing
from the probability distribution.

The matrix is nonsingular. Each state in $U$ has a simple positive path
to $t$ of length at most $r$. The minimum, over finitely many chosen
paths, of their probability products is some $\delta>0$. Therefore
from every state of $U$ the probability of still being in $U$ after
$r$ steps is at most $1-\delta$. Hence
$\|Q^r\|_\infty\leq1-\delta$, the Neumann series
$\sum_{h\geq0}Q^h$ converges, and it is the inverse of $I-Q$.
This also verifies that the solution of \eqref{a478:linear} is the
probabilistic one. The case $U=\varnothing$ requires no linear solve.

Clear denominators row by row and use exact rational or fraction-free
Gaussian elimination. The next lemma bounds the relevant minors by
integers of polynomial bit length, so this is a polynomial-time bit
algorithm, not merely a count of unit-cost operations on unbounded
rationals. States outside $R$ were removed first; trying to invert the
full matrix with an unreachable closed class would be incorrect.
\hfill$\square$

\paragraph{Lemma 478.3: a uniform denominator bound.}
Let $m$ be the number of nonterminal vertices in the original game. Let
$L$ be the sum of binary encoding lengths of all positive chance-transition
denominators, including repeated occurrences in different rows. Put
\begin{equation}\label{a478:denominator}
 B=L+m\lceil\log_2(m+1)\rceil+1,\qquad D=2^B.
\end{equation}
Every fixed-policy value, and every game value, has a reduced rational
denominator at most $D$. Numerators of values in $[0,1]$ are bounded by
the same number. In particular exact output length is polynomial in the
input length.

\emph{Proof for a fixed pair.} In the system from Lemma 478.2 multiply
row $i$ by $d_i$, the product of the denominators in its transition row;
deterministic rows use $d_i=1$. Every resulting matrix entry has absolute
value at most $d_i$, and the right-hand entry also has absolute value at
most $d_i$, because all the underlying probabilities are in $[0,1]$.
If the system has $r\leq m$ rows, Hadamard's determinant bound gives
\[
       1\leq|\det M|\leq r^{r/2}\prod_i d_i
                     \leq m^{m/2}2^L\leq D
\]
for $r>0$, with trivial interpretation when $m=0$. Cramer's rule says
that each reduced denominator divides the nonzero integer $\det M$.
Replacing one column by the right-hand side gives the same bound for
the numerator determinant. The bound also applies to minors encountered
in fraction-free elimination, with smaller dimensions, proving the bit
complexity assertion in the previous lemma.

For the game-value conclusion, use the established deterministic
memoryless optimality theorem for finite turn-based stochastic reachability
games, stated in [E2], Section 2.2. Its hypotheses match this graph:
finite state and action sets, observed turns, reachability objective, and
fixed transition probabilities. It does not require all policy pairs
to terminate. It supplies optimal stationary policies whose induced
chain's hitting probabilities equal the game values. Applying the
fixed-pair bound proves the assertion. Existence of those policies is an
external theorem here, not an algorithm for discovering them.
\hfill$\square$

The determinant estimate deliberately uses a generous bound. It is enough
that $B$ is polynomial in the actual rational input length; no claim of
a strongly polynomial algorithm independent of probability bit lengths
is being made. Large numerical denominators can occur with short binary
encodings, and their \emph{bit lengths}, rather than their values, belong
in the complexity accounting.

\paragraph{Exact thresholds and approximation are linked only at controlled precision.}
Two distinct reduced rationals in $[0,1]$ with denominators at most $D$
differ by at least $D^{-2}$, since their cross-multiplied numerators differ
by a nonzero integer. Consequently an oracle deciding $v(s)\geq q$
can isolate a value in an interval of width less than $1/(2D^2)$ with
$O(B)$ binary-search queries. The dyadic thresholds have $O(B)$ bits.
If the oracle answers yes at a midpoint, retain that midpoint as the
lower endpoint; otherwise retain it as the upper endpoint. Equality is
handled by the yes branch, rather than being replaced by a strict test.

The unique rational with denominator at most $D$ in the resulting interval
can be recovered by continued fractions in polynomial bit complexity.
For an explicit criterion, take its midpoint $z$, so
$|z-v|<1/(4D^2)\leq1/(2\operatorname{den}(v)^2)$. The elementary
continued-fraction approximation criterion then makes $v$ a convergent
of $z$; list its convergents and choose the one inside the isolating
interval with denominator at most $D$. The criterion follows from the
best-approximation property of consecutive convergents: if a different
$a/b$ were that close, the nonzero integer difference of the cross
products with the last convergent of denominator at most $b$ would have
absolute value less than one. The Euclidean algorithm uses polynomially
many bit operations on the $O(B)$-bit endpoints. Values zero and one
are included, or may be checked at the outset.

Conversely an exact rational value immediately decides the threshold.
For a threshold $q=a/b$ in lowest terms, a nonequal value differs from
$q$ by at least $1/(Db)$. But an approximation alone cannot decide
what to do in the equality case; rational reconstruction or a certified
isolating procedure is needed. For example a fair chance between target
and failure has $v=1/2$, so the statements $v\geq1/2$ and $v>1/2$
have opposite answers. This distinction is not removed by saying that
probabilities are approximated accurately.

A uniform approximation algorithm polynomial in the input length and
$\log(1/\varepsilon)$, with a certified error bound, would suffice:
use $\varepsilon<1/(4D^2)$ and reconstruct the rational. A bound
polynomial in $1/\varepsilon$ does not imply polynomial time at that
precision. The denominator lemma identifies exactly why the two precision
dependencies are different for this question.

\paragraph{A concrete obstruction to direct value iteration.}
Take one chance vertex with a self-loop of probability $1-2^{-b}$ and
a transition to target of probability $2^{-b}$. Its rational encoding
uses $O(b)$ bits. The exact value is one: the probability of avoiding
the target for $h$ steps is $(1-2^{-b})^h$, which tends to zero.
The iterates in Lemma 478.1 are exactly
\begin{equation}\label{a478:rare}
                        x^{(h)}=1-(1-2^{-b})^h.
\end{equation}
For $b\geq2$ and $h\leq2^{b-1}$, the union bound for success in
one of $h$ trials gives $x^{(h)}\leq h2^{-b}\leq1/2$;
at $h=2^{b-1}>1$ the inequality is strict because the success events
overlap in the independent-trial representation. Thus exponentially many
iterations may be needed even to exceed the constant threshold $1/2$.
This chain is absorbing with probability one, so merely excluding
nonstopping games does not repair the iteration count.

The same issue occurs with only fair probabilities. Use states
$0,1,\ldots,r$, with $r$ the target. From $i<r$, a fair head moves
to $i+1$ and a fair tail resets to $0$. Hitting the target means observing
$r$ consecutive heads. Its value is one: successive disjoint blocks of
$r$ tosses independently have probability $2^{-r}$ of all heads, so
the probability that no such block appears tends to zero. By time $h$
there are at most $h$ candidate ending positions for a run, hence
\begin{equation}\label{a478:fair}
                    x^{(h)}_0\leq h2^{-r}.
\end{equation}
At $h=2^{r-1}$ this is at most $1/2$, despite all transition probabilities
having constant bit length and the graph having only $r+1$ vertices.
Lemma 478.2 solves this chain in polynomial time. The example defeats
one iterative method, not the existence of some different polynomial
algorithm for games.

\paragraph{A complete one-player linear-programming benchmark.}
If there are no Min vertices, consider the rational linear program
\begin{equation}\label{a478:LP}
 \begin{gathered}
 \text{minimize }\sum_{i\ne t}z_i,\qquad 0\leq z_i\leq1,
 \quad z_t=1,\\
 z_i\geq z_j\quad(i\text{ a Max vertex},\ j\in\operatorname{Succ}(i)),\\
 z_i\geq\sum_jp_{ij}z_j\quad(i\text{ a chance vertex}),
 \end{gathered}
\end{equation}
with failure terminals fixed at zero. Every feasible vector satisfies
$z\geq Tz$ and dominates the initial vector, so monotonicity gives
$z\geq x^{(h)}$ for all $h$, hence $z\geq v$. The value vector is
itself feasible by Lemma 478.1. Therefore it is the unique coordinatewise
minimal feasible vector and the unique minimizing vector. There are
polynomially many variables and constraints of polynomial rational bit
length. The standard polynomial-time theorem for rational linear
programming makes this a polynomial bit algorithm for this one-player
subcase, including nonstopping components.

Naively dualizing this program for Min is dangerous. At a Min vertex
with self-loop and target actions, the inequalities $z\leq z$ and
$z\leq1$, together with maximizing $z$, would choose one, whereas
the actual value is zero. Correct one-player Min algorithms handle the
zero-value region and recurrent components; the example shows why the
unqualified reversed program is wrong.

For two players, merely combining the Max lower inequalities with the
Min upper inequalities does not select the Bellman solution. Even the
acyclic chain $M\to N\to t$, with $M$ a Max vertex and $N$ a Min
vertex and each having just one action, has true value one at both
vertices. The mixed inequalities only say $z_M\geq z_N$ and
$z_N\leq1$, and permit $z_M=z_N=0$. Conversely the deterministic
chain from a Min vertex through a Max vertex to a failure terminal admits
spurious positive vectors under the opposite objective. Local inequality
feasibility is not a general certificate of the correct strategic value.

There is an additional policy-selection trap even with the correct vector
already in hand. At a Max vertex with self-loop and target actions the
value is one, and both actions lead to a state of value one. Choosing
the self-loop forever is locally maximizing but has success probability
zero. The stationary optimality theorem guarantees a suitable choice,
not that every arbitrary greedy tie choice is suitable. End-component
control is essential.

\paragraph{Exact finite verification.}
The companion script exhaustively examines 3,375 five-vertex games with
one Max vertex, one Min vertex, one fair-chance vertex, target, and failure.
At each nonterminal vertex its two successor positions range over all
15 unordered pairs with repetition from the five states. Self-loops and
unreachable closed components are included. For every game the four
stationary policy pairs are evaluated using rational Gaussian elimination
after removing states with no positive path to target: 13,500 exact
fixed-policy evaluations in total.

For each game, componentwise max--min and min--max values agree, there
is a common optimal stationary strategy for each player across all states,
the Bellman equations hold, and the zero-region computation
\eqref{a478:zero} agrees with the exact zero coordinates. The first eight
finite-horizon iterates are checked to lie below those values. All
denominators satisfy the conservative small-instance bound used by the
script. These finite checks test the implementation and quantifier order;
they do not replace the cited general stationary-optimality theorem.

The rare-transition examples are checked exactly for $b=2,4,8,12$,
including their values after $2^{b-1}$ steps. The fair-run chains are
checked for $2\leq r\leq10$ after $2^{r-1}$ steps. In each case the
exact infinite-horizon value is one while the iteration remains at most
one half. No floating-point tolerance is used. Scripts, complete results,
and source/review records are saved under
\nolinkurl{_Local/top_problems/continuation_488_477_2026_09_27/algorithms/}.

\paragraph{Outcome and first remaining gap.}
\textbf{UNRESOLVED.} This attempt establishes the least-fixed-point
characterization, a polynomial qualitative zero test, exact fixed-policy
evaluation, an explicit polynomial bound on value bit lengths, and the
precision needed to recover exact rational values. It proves a one-player
linear-programming subcase and supplies two exponential-iteration families
and several counterexamples to proposed local shortcuts.

The missing step is a polynomial-time method to select or implicitly
represent optimal choices for both players while retaining the least
reachability solution. Enumerating all stationary policies is finite and
correct using the cited optimality theorem, but the number of policy pairs
can be exponential in the explicit graph size. Polynomial cost per pair
does not make that enumeration polynomial. Nor does the short rational
output supply a fast way to find it.

Specific next targets are a globally valid potential bounding strategy
changes, a polynomial-size convex formulation that enforces recurrence
conditions, or an approximation method with polynomial dependence on
$\log(1/\varepsilon)$. Each must be tested on the fair-run chain,
the Min self-loop example, and games with both kinds of decisions.
No such global bound or formulation is proved here, and neither the
exact-value problem nor its non-strict threshold version is resolved.

\subsection{Sources}
\begin{itemize}
\item[S1]Formal statement, dblp record header for journals/iandc/Condon92.
\url{https://dblp.org/rec/journals/iandc/Condon92.html}.
\textit{Formulation source.}
\item[E1]Anne Condon, The Complexity of Stochastic Games,
Information and Computation 96 (1992), 203--224.
\url{https://www.sciencedirect.com/science/article/pii/089054019290048K}.
\textit{Primary publication corresponding to [S1]; abstract and publication
record consulted 27 September 2026.}
\item[E2]Jan K\v{r}et\'{\i}nsk\'y, Emanuel Ramneantu, Alexander Slivinskiy,
and Maximilian Weininger, Comparison of Algorithms for Simple Stochastic
Games (Full Version), arXiv:2008.09465v3, 25 August 2020,
Sections 2 and 3, and the discussion of end components.
\url{https://arxiv.org/pdf/2008.09465}.
\textit{Primary algorithm paper covering memoryless optimality, Bellman
semantics, and nonstopping-game issues; consulted 27 September 2026.}
\item[E3]Avi Rudich, Isaac Rudich, and Rachel Rue, Simple Stochastic
Stopping Games: A Generator and Benchmark Library,
arXiv:2402.02571 (2024).
\url{https://arxiv.org/abs/2402.02571}.
\textit{Primary benchmark paper's abstract used only for context on the
unresolved algorithmic question, not as a worst-case algorithm theorem.
Consulted 27 September 2026.}
\end{itemize}
\end{document}
